\documentclass[10pt,a4paper]{amsart}
\usepackage[margin=19mm]{geometry}
\usepackage{amsmath,amssymb,mathtools}
\usepackage[scr=rsfs,cal=euler]{mathalfa}
\usepackage{xfrac}
\usepackage[unicode,hidelinks,bookmarks=false]{hyperref}
\newcommand{\ie}{i.e.\ }
\newcommand{\cf}{cf.\ }
\newcommand{\resp}{respectively\ }
\newcommand{\Subsection}{\subsection}
\providecommand{\nobreakdash}{\nobreak\hbox{-}}
\setlength{\emergencystretch}{2em}
\multlinegap0pt
\usepackage{bm}
\usepackage{bbm}
\usepackage{dsfont}
\usepackage{xspace}
\usepackage{cancel}
\usepackage{mathtools}
\usepackage{amsthm}
\makeatletter
\@ifundefined{thm}{\newtheorem{thm}{Theorem}}{}
\@ifundefined{rem}{\newtheorem{rem}[thm]{Remark}}{}
\makeatother
\title[Some interpolation inequalities in Lorentz, Morrey and BMO s]{Some interpolation inequalities in Lorentz, Morrey and BMO spaces}
\author{Hua Wang}
\date{Source: arXiv:2511.06589v1 (CC BY 4.0)}
\begin{document}
\maketitle

\noindent\emph{Source and licence.} Some interpolation inequalities in Lorentz, Morrey and BMO spaces. Version arXiv:2511.06589v1 (CC BY 4.0), distributed under the Creative Commons Attribution 4.0 International licence, as recorded in the arXiv licence metadata for this exact version and shown on the abstract page https://arxiv.org/abs/2511.06589v1. Full rights notice: licenses/arxiv-2511.06589-CC-BY-4.0.txt.\par\smallskip

\noindent\textit{Editorial scope.} Counted unit: the Rem whose source label is \texttt{22} in the article (source0.tex lines 584--590, source label \texttt{22}), delivered together with the article's own complete original proof of it (source0.tex lines 592--756, one \texttt{proof} environment of 165 lines). No mathematics is written, repaired or re-derived: statement and proof are the article's own text line for line. Required context retained uncounted: equality2 (lines 146--152); geq (lines 178--180); weneed (lines 462--465); thmwang1 (lines 475--490); strong (lines 492--498). Every definition, display and label of those lines is retained unchanged. Omitted from this package and not redistributed: 1--145, 153--177, 181--461, 466--474, 491--491, 499--583, 591--591, 757--1533. Nothing inside a retained excerpt is omitted or altered: every retained body is delivered exactly as the article's lines are written, so no omission receipt of any kind accompanies this package. Citations: the retained excerpts carry no citation command at all, so no bibliography entry of the article is transcribed and no reference is redistributed. Reference adaptation: none was needed and none was made; every reference inside the delivered unit resolves inside it, so no unresolved reference or citation is produced. Printed slips kept literal: no printed word, symbol, hypothesis or number of the counted statement or of its proof was altered. Layout and class: the article's own class is \texttt{birkjour} with packages including amsmath, amssymb; the wrapper is the lane's pinned \texttt{amsart} editorial wrapper at 10pt on A4, which restates the theorem-like environments the retained text needs and adds the article's own macro definitions that the retained text needs (none), with extra wrapper packages (bm, bbm, dsfont, xspace, cancel, mathtools). The delivered numbering is the unit's own. Assets: no retained span contains a figure environment, a graphics-inclusion command, a PDF-page inclusion, a table or a TikZ picture; no image file is copied into this package, the package declares no asset dependency and no image file is redistributed with it. Licence: version arXiv:2511.06589v1 of this article is distributed under the Creative Commons Attribution 4.0 International licence, as recorded in the arXiv licence metadata for this exact version and shown on the abstract page https://arxiv.org/abs/2511.06589v1. A full rights notice is delivered at licenses/arxiv-2511.06589-CC-BY-4.0.txt. Characters: every retained excerpt is pure ASCII, so the delivered engine, XeLaTeX with its default Latin Modern text font, reports no missing character.
\begin{equation}\label{equality2}
\|f\|_{L^p}=
\begin{cases}
\displaystyle\bigg(\int_{0}^{\infty}\big[f^{*}(t)\big]^p\,dt\bigg)^{1/p}, &\mbox{if}~ p\in[1,\infty),\\
\displaystyle\sup_{t>0}f^{*}(t)=\lim_{t\downarrow 0}f^{*}(t), &\mbox{if}~ p=\infty.
\end{cases}
\end{equation}

\begin{equation}\label{geq}
f^{**}(t)-f^{*}(t)\geq 0,\quad t>0.
\end{equation}

\begin{equation}\label{weneed}
\bigg(\int_0^{t}\Big[\log\frac{t}{\,s\,}\Big]^qds\bigg)^{1/q}
=\bigg(\int_0^{\infty}t\cdot\mu^q e^{-\mu}\,d\mu\bigg)^{1/q}=t^{1/q}\cdot\Gamma(q+1)^{1/q},
\end{equation}

\begin{thm}\label{thmwang1}
Let $n\in \mathbb{N}$, $1\leq p<\infty$ and $p<r\leq\infty$. Then the following statements hold.

$(1)$ If $r<\infty$ and $f\in L^{p,r}(\mathbb R^n)\cap{L^\infty}(\mathbb R^n)$, then for every $q$ with $r<q<\infty$, we have $f\in L^q(\mathbb R^n)$, and the interpolation inequality
\begin{equation*}
\|f\|_{L^q}\leq 2\cdot\big(\|f\|_{L^{p,r}}\big)^{p/q}\cdot\big(\|f\|_{L^\infty}\big)^{1-p/q}
\end{equation*}
holds for all $f\in L^{p,r}(\mathbb R^n)\cap{L^\infty}(\mathbb R^n)$.

$(2)$ If $r=\infty$ and $f\in L^{p,r}(\mathbb R^n)\cap{L^\infty}(\mathbb R^n)$, then for every $q$ with $p<q<\infty$, we have $f\in L^q(\mathbb R^n)$, and the interpolation inequality
\begin{equation*}
\|f\|_{L^q}\leq \Big[2^{1-1/q}\Big(\frac{q}{q-p}\Big)^{1/q}\Big]
\cdot\big(\|f\|_{L^{p,\infty}}\big)^{p/q}\cdot\big(\|f\|_{L^\infty}\big)^{1-p/q}
\end{equation*}
holds for all $f\in L^{p,\infty}(\mathbb R^n)\cap{L^\infty}(\mathbb R^n)$.
\end{thm}

\begin{thm}\label{strong}
Let $n\in \mathbb{N}$, $1<p<\infty$ and $p\leq r\leq\infty$. If $p\leq r<\infty$ and $f\in L^{p,r}(\mathbb R^n)\cap \mathcal{W}$, then for every $q$ with $r<q<\infty$(or if $f\in L^{p,\infty}(\mathbb R^n)\cap \mathcal{W}$, then for every $q$ with $p<q<\infty$), we have $f\in L^q(\mathbb R^n)$, and the following interpolation inequality
\begin{equation*}
\|f\|_{L^q}\leq C\cdot q\big(\|f\|_{L^{p,r}}\big)^{p/q}\cdot\big(\|f\|_{\mathcal{W}}\big)^{1-p/q}
\end{equation*}
holds for all $f\in L^{p,r}(\mathbb R^n)\cap \mathcal{W}$. Here $C$ is an absolute positive constant independent of $q$.
\end{thm}

\begin{rem}
For $1\leq p<\infty$, it is easy to verify that
\begin{equation}\label{22}
\lim_{q\to\infty}2^{1-1/q}\Big(\frac{q}{q-p}\Big)^{1/q}=2,
\end{equation}
that is to say $2^{1-1/q}\big(\frac{q}{q-p}\big)^{1/q}=\mathcal{O}(2)$, as $q\to\infty$. Theorem $\ref{thmwang1}$ tells us that the embedding constant for the case $p<r<\infty$ is 2 and the asymptotically embedding constant for the case $p<r=\infty$ is also 2.
\end{rem}

\begin{proof}[Proof of Theorem $\ref{strong}$]
Let $f\in L^{p,r}(\mathbb R^n)\cap \mathcal{W}$ with $1<p<\infty$ and $p\leq r\leq\infty$, and let $p'$ be the exponent conjugate to $p$. First we claim that for any $t>0$,
\begin{equation}\label{key1}
f^{**}(t)\leq\Big(\frac{p'}{r'}\Big)^{1/{r'}}
\frac{1}{t^{1/p}}\cdot\|f\|_{L^{p,r}}.
\end{equation}
Here $r'=1$ when $r=\infty$.
In fact, we first suppose $p\leq r<\infty$. For any $t>0$, by using H\"{o}lder's inequality with exponent $r$, we obtain
\begin{equation*}
\begin{split}
\int_0^{t}f^{*}(\tau)\,d\tau&=\int_0^{t}\tau^{1/p-1/r}f^{*}(\tau)\cdot\tau^{1/r-1/p}\,d\tau\\
&\leq\bigg(\int_0^{t}\Big[\tau^{1/p}f^{*}(\tau)\Big]^r\frac{d\tau}{\tau}\bigg)^{1/r}
\times\bigg(\int_0^{t}\tau^{(1/r-1/p)r'}d\tau\bigg)^{1/{r'}}\\
&\leq \|f\|_{L^{p,r}}\bigg(\int_0^{t}\tau^{(1/r-1/p)r'}d\tau\bigg)^{1/{r'}}.
\end{split}
\end{equation*}
Recall that $p>1$. A simple computation leads to that
\begin{equation*}
(1/r-1/p)r'+1>0.
\end{equation*}
Then we have
\begin{equation*}
\begin{split}
\bigg(\int_0^{t}\tau^{(1/r-1/p)r'}\,d\tau\bigg)^{1/{r'}}
&=\Big[\frac{1}{(1/r-1/p)r'+1}\cdot t^{(1/r-1/p)r'+1}\Big]^{1/{r'}}\\
&=\Big(\frac{1}{r'}\Big)^{1/{r'}}\Big(\frac{1}{1-1/p}\Big)^{1/{r'}}\cdot t^{1-1/p}.
\end{split}
\end{equation*}
Hence, for any $t>0$,
\begin{equation*}
\begin{split}
f^{**}(t)&=\frac{\,1\,}{t}\int_0^{t}f^{*}(\tau)\,d\tau\\
&\leq\frac{\,1\,}{t}\Big(\frac{1}{r'}\Big)^{1/{r'}}\big(p'\big)^{1/{r'}}
\cdot t^{1-1/p}\|f\|_{L^{p,r}}\\
&=\Big(\frac{p'}{r'}\Big)^{1/{r'}}\frac{1}{t^{1/p}}\cdot\|f\|_{L^{p,r}},
\end{split}
\end{equation*}
as desired. Now we suppose $r=\infty$. In this case, it is easy to verify that
\begin{equation*}
\begin{split}
\int_0^{t}f^{*}(\tau)\,d\tau&=\int_0^{t}\tau^{1/p}f^{*}(\tau)\cdot\tau^{-1/p}\,{d\tau}\\
&\leq\Big[\sup_{\tau>0}\tau^{1/p}f^{*}(\tau)\Big]\cdot\int_0^{t}\tau^{-1/p}\,{d\tau}\\
&=\|f\|_{L^{p,\infty}}\frac{p}{p-1}\cdot t^{1-1/p},
\end{split}
\end{equation*}
where in the last step we have used the fact that $p>1$. Therefore, for any $t>0$,
\begin{equation*}
\begin{split}
f^{**}(t)&=\frac{\,1\,}{t}\int_0^{t}f^{*}(\tau)\,d\tau\\
&\leq p'\frac{1}{t^{1/p}}\cdot\|f\|_{L^{p,\infty}}.
\end{split}
\end{equation*}
This proves \eqref{key1}. It should be pointed out that inequality \eqref{key1} fails when $p=1$. On the other hand, since
\begin{equation}\label{important}
\begin{split}
\frac{d}{d\tau}\big(f^{**}(\tau)\big)&=-\frac{1}{\tau^2}\int_0^{\tau}f^{*}(\eta)\,d\eta+\frac{f^{*}(\tau)}{\tau}\\
&=\frac{f^{*}(\tau)-1/{\tau}\int_0^{\tau}f^{*}(\eta)\,d\eta}{\tau}\\
&=\frac{f^{*}(\tau)-f^{**}(\tau)}{\tau},
\end{split}
\end{equation}
so we have
\begin{equation*}
-\frac{d}{d\tau}\big(f^{**}(\tau)\big)=\frac{f^{**}(\tau)-f^{*}(\tau)}{\tau}\leq\frac{\|f\|_{\mathcal{W}}}{\tau},
\end{equation*}
for given $f\in \mathcal{W}$. From this, it follows that for any $0<s<t$,
\begin{equation*}
\begin{split}
f^{**}(s)-f^{**}(t)&=\int_{s}^t-\frac{d}{d\tau}\big(f^{**}(\tau)\big)\,d\tau\\
&\leq\int_{s}^t\frac{\|f\|_{\mathcal{W}}}{\tau}\,d\tau\\
&=\|f\|_{\mathcal{W}}\cdot\big[\log t-\log s\big]=\|f\|_{\mathcal{W}}\cdot\log\frac{t}{\,s\,},
\end{split}
\end{equation*}
which in turn implies that
\begin{equation}\label{key2}
f^{**}(s)\leq f^{**}(t)+\|f\|_{\mathcal{W}}\cdot\log\frac{t}{\,s\,},
\end{equation}
whenever $0<s<t$. For any $1<q<\infty$, in view of \eqref{equality2}, we can write
\begin{equation*}
\begin{split}
\|f\|_{L^q}&=\bigg(\int_0^{\infty}\big[f^{*}(s)\big]^qds\bigg)^{1/q}\\
&\leq\bigg(\int_0^{t}\big[f^{*}(s)\big]^qds\bigg)^{1/q}+
\bigg(\int_{t}^{\infty}\big[f^{*}(s)\big]^qds\bigg)^{1/q}\\
&:=J_1(t)+J_2(t),
\end{split}
\end{equation*}
where $t$ is a positive constant to be fixed below. Let us estimate the first term $J_1(t)$. By using \eqref{geq}, \eqref{key2} and Minkowski's inequality, we can deduce that
\begin{equation*}
\begin{split}
J_1(t)&\leq\bigg(\int_0^{t}\big[f^{**}(s)\big]^qds\bigg)^{1/q}\\
&\leq\bigg(\int_0^{t}\big[f^{**}(t)\big]^qds\bigg)^{1/q}
+\bigg(\int_0^{t}\Big[\|f\|_{\mathcal{W}}\cdot\log\frac{t}{\,s\,}\Big]^qds\bigg)^{1/q}\\
&=t^{1/q}\cdot f^{**}(t)+\|f\|_{\mathcal{W}}
\cdot\bigg(\int_0^{t}\Big[\log\frac{t}{\,s\,}\Big]^qds\bigg)^{1/q}.\\
\end{split}
\end{equation*}
Recall that the Stirling formula
\begin{equation*}
\Gamma(q+1)\approx\sqrt{2\pi}q^{q+\frac{\,1\,}{2}}\exp(-q)=\sqrt{2\pi q}\Big(\frac{q}{\,e\,}\Big)^q
\end{equation*}
holds as $q\to\infty$. Hence, we find that as $q\to\infty$,
\begin{equation}\label{aso}
\Gamma(q+1)^{1/q}=\mathcal{O}(q),
\end{equation}
that is, $\Gamma(q+1)^{1/q}\leq C\cdot q$, when $q$ is large. Here $C>0$ is an absolute constant independent of $q$. This fact, together with \eqref{weneed} and \eqref{key1}, gives us that
\begin{equation*}
\begin{split}
J_1(t)&\leq C\cdot q\Big[t^{1/q-1/p}\cdot\|f\|_{L^{p,r}}+\|f\|_{\mathcal{W}}\cdot t^{1/q}\Big],
\end{split}
\end{equation*}
where $C>0$ is independent of $f$ and $q>1$. We now estimate the second term $J_2(t)$.

$(1)$ When $r<\infty$ and $p\leq r<q<\infty$, since $f^{*}$ is non-increasing and $r/p-1\geq0$, we have
\begin{equation*}
\begin{split}
J_2(t)&=\bigg(\int_{t}^{\infty}\big[s^{1/p-1/r}f^{*}(s)\big]^r\cdot f^{*}(s)^{q-r}s^{1-r/p}ds\bigg)^{1/q}\\
&\leq\bigg(\int_{t}^{\infty}\big[s^{1/p-1/r}f^{*}(s)\big]^rds\bigg)^{1/q}\cdot \Big[f^{*}(t)^{q-r}t^{1-r/p}\Big]^{1/q}\\
&\leq\bigg(\int_{0}^{\infty}\big[s^{1/p}f^{*}(s)\big]^r\frac{ds}{s}\bigg)^{1/q}\cdot \Big[f^{*}(t)^{q-r}t^{1-r/p}\Big]^{1/q}.
\end{split}
\end{equation*}
Moreover, by \eqref{geq}, \eqref{key1} and the fact that $q-r>0$, we thus obtain
\begin{equation*}
\begin{split}
J_2(t)&\leq\big(\|f\|_{L^{p,r}}\big)^{r/q}
\cdot \Big[f^{**}(t)^{q-r}t^{1-r/p}\Big]^{1/q}\\
&\leq \big(\|f\|_{L^{p,r}}\big)^{r/q}
\cdot \Big[\Big(\frac{p'}{r'}\Big)^{1/{r'}}\Big]^{1-r/q}
t^{1/q-1/p}\big(\|f\|_{L^{p,r}}\big)^{1-r/q}\\
&=\Big[\Big(\frac{p'}{r'}\Big)^{1/{r'}}\Big]^{1-r/q}
t^{1/q-1/p}\cdot\|f\|_{L^{p,r}}.
\end{split}
\end{equation*}
$(2)$ When $r=\infty$ and $p<q<\infty$, in this case, we have
\begin{equation*}
\begin{split}
J_2(t)&=\bigg(\int_{t}^{\infty}\big[s^{1/p}f^{*}(s)\big]^{q}\frac{1}{s^{q/p}}ds\bigg)^{1/q}\\
&\leq\Big[\sup_{s>0}s^{1/p}f^{*}(s)\Big]\cdot\bigg(\int_{t}^{\infty}\frac{1}{s^{q/p}}ds\bigg)^{1/q}\\
&=\Big(\frac{p}{q-p}\Big)^{1/q} t^{1/q-1/p}\cdot\|f\|_{L^{p,\infty}},
\end{split}
\end{equation*}
where the last integral is convergent since $q/p>1$. Notice that $p'\geq r'$. Thus, for all $q$ with $p<q<\infty$,
\begin{equation*}
\Big[\Big(\frac{p'}{r'}\Big)^{1/{r'}}\Big]^{1-r/q}\leq\Big(\frac{p'}{r'}\Big)^{1/{r'}}.
\end{equation*}
Moreover, it is easy to see that
\begin{equation*}
\lim_{q\to\infty}\Big(\frac{p}{q-p}\Big)^{1/q}=1.\Longrightarrow \Big(\frac{p}{q-p}\Big)^{1/q}=\mathcal{O}(1)\quad \mbox{as}~ q\to\infty.
\end{equation*}
Summing up the above estimates, we then conclude that
\begin{equation*}
\|f\|_{L^q}\lesssim q\Big[t^{1/q-1/p}\cdot\|f\|_{L^{p,r}}+\|f\|_{\mathcal{W}}\cdot t^{1/q}\Big]
\end{equation*}
with the implicit positive constant independent of $q$ (may depend on $p$ and $r$). We now take $t>0$ in this estimate such that
\begin{equation*}
t^{1/q-1/p}\cdot\|f\|_{L^{p,r}}=\|f\|_{\mathcal{W}}\cdot t^{1/q}.
\end{equation*}
In other words,
\begin{equation*}
t:=\left(\frac{\|f\|_{L^{p,r}}}{\|f\|_{\mathcal{W}}}\right)^p.
\end{equation*}
Therefore, by such choice of $t>0$, we finally obtain
\begin{equation*}
\|f\|_{L^q}\lesssim q\cdot\big(\|f\|_{L^{p,r}}\big)^{p/q}\cdot\big(\|f\|_{\mathcal{W}}\big)^{1-p/q},
\end{equation*}
where the implicit positive constant is independent of $q>1$. We are done.
\end{proof}

\end{document}
